\documentclass[11pt,a4paper]{article}

\usepackage[utf8]{inputenc}
\usepackage[T1]{fontenc}
\usepackage[english]{babel}
\usepackage{geometry}
\geometry{a4paper, margin=2.5cm}

\usepackage{booktabs}
\usepackage{array}
\usepackage{multirow}
\usepackage{tabularx}
\usepackage{longtable}
\usepackage{hyperref}
\usepackage{url}
\usepackage{amsmath}
\usepackage{amssymb}
\usepackage{graphicx}
\usepackage{xcolor}
\usepackage{enumitem}
\usepackage{caption}
\usepackage{subcaption}
\usepackage{parskip}
\usepackage{microtype}
\usepackage{fancyhdr}
\usepackage{titlesec}
\usepackage[backend=bibtex,style=ieee,sorting=none]{biblatex}

\hypersetup{
    colorlinks=true,
    linkcolor=blue!70!black,
    citecolor=blue!70!black,
    urlcolor=blue!60!black,
}

\pagestyle{fancy}
\fancyhf{}
\rhead{Try 66 — State-of-the-Art Comparison}
\lhead{TFG Path Loss Prediction}
\rfoot{\thepage}

\definecolor{highlight}{RGB}{0,100,180}
\definecolor{good}{RGB}{0,140,0}
\definecolor{caution}{RGB}{180,90,0}

\newcommand{\better}[1]{\textcolor{good}{\textbf{#1}}}
\newcommand{\ours}[1]{\textcolor{highlight}{\textbf{#1}}}
\newcommand{\yes}{$\checkmark$}
\newcommand{\no}{$\times$}

\title{%
    \textbf{State-of-the-Art Comparison for Outdoor Path Loss Prediction}\\[0.5em]
    \large PMHHNet with Sinusoidal FiLM Conditioning (Try 66)\\[0.3em]
    \normalsize TFG — Grau en Enginyeria en Telecomunicació
}
\author{Guillem Moreno}
\date{April 2026}

\begin{document}

\maketitle
\thispagestyle{fancy}

\begin{abstract}
This document compares the path loss prediction performance of our model (PMHHNet,
herein referred to as Try 66) against published state-of-the-art methods for
outdoor and indoor radio map estimation. The key distinguishing features of our
work are: (i) city-level holdout evaluation over 66 OSM-derived city geometries
with simulator-generated propagation labels,
(ii) continuous UAV transmitter height conditioning via sinusoidal Feature-wise
Linear Modulation (FiLM), and (iii) 513$\times$513 pixel spatial resolution.
We show that Try 66 achieves 9.41~dB root mean square error (RMSE) on the
\texttt{open\_sparse\_lowrise} expert, which is competitive with the best published
results for cross-environment generalization and numerically equals the ICASSP 2025
indoor challenge winner, despite our problem being structurally harder in every
measured dimension.
\end{abstract}

\tableofcontents
\newpage

%---------------------------------------------------------------------
\section{Introduction}
\label{sec:intro}
%---------------------------------------------------------------------

Radio map estimation — predicting path loss at every spatial location given a
transmitter position and an environment map — is a core problem in wireless
network planning, digital twins, and UAV-assisted communications. Deep learning
approaches have rapidly replaced classical empirical and ray-tracing methods for
this task, achieving sub-2~dB RMSE on fixed-Tx benchmark datasets
\cite{pmnet2023,remnet2024,rmtransformer2025}.

However, the dominant benchmarks (primarily the ICASSP 2023 Outdoor Pathloss
Challenge \cite{icassp2023challenge}) share a critical limitation: both training
and test environments are generated by the \emph{same} procedural city simulator,
and the transmitter is always fixed at a rooftop height of 6--20~m. These
constraints allow models to implicitly memorize the generator's statistical patterns
rather than learn true physical generalisation.

This document contextualises our results within the broader literature by:
\begin{enumerate}[nosep]
    \item Examining all known outdoor and indoor benchmarks with their key conditions.
    \item Identifying which comparisons are methodologically valid.
    \item Highlighting the genuinely novel aspects of Try 66 that no current paper
          addresses simultaneously.
\end{enumerate}

%---------------------------------------------------------------------
\section{Problem Formulation}
\label{sec:problem}
%---------------------------------------------------------------------

\subsection{Try 66 Setup}

\textbf{Dataset:} \texttt{CKM\_Dataset\_270326.h5} — 16,180 samples spanning 66
OSM-derived city geometries, stored as \texttt{uint8} quantised path loss values (1~dB
resolution). Each sample is a 513$\times$513 pixel map at 7.125~GHz.

\textbf{Inputs:} Eight channels per sample: topology map, LoS mask, distance map,
formula prior (free-space + correction), confidence mask, Tx depth map, elevation
angle map, and building mask.

\textbf{Evaluation split:} City-level holdout over unseen city geometries —
validation cities are entirely excluded from training. No fine-tuning or
calibration at test time.

\textbf{Transmitter height:} Continuous UAV altitude in range 12--478~m, conditioned
via sinusoidal positional encoding fed into every FiLM layer.

\textbf{Architecture:} PMHHNet — a PMNet backbone augmented with a high-frequency
(HF) branch processing Laplacian-filtered inputs and FiLM conditioning on antenna
height. Squeeze-and-Excitation (SE) attention in the encoder.

\textbf{Training:} AdamW optimizer, Huber loss with NLoS reweighting
($w_\mathrm{NLoS} = 6\times$), multiscale loss, CutMix augmentation, cosine warm
restarts (SGDR), EMA weight averaging.

\subsection{Metrics and Conventions}

All RMSE values in this document are in decibels (dB) unless otherwise noted.
Where papers report normalised RMSE for the ICASSP 2023 outdoor challenge, we
convert to dB using the actual RadioMap3DSeer image mapping reported in the
dataset paper \cite{dataset2212}. For the rooftop 3D dataset used by the challenge,
the truncation threshold and maximum value are
\[
    \mathrm{PL}_{\mathrm{thr}} = -104~\mathrm{dB}, \qquad
    \mathrm{PL}_{\max} = -75~\mathrm{dB},
\]
and the image-domain normalisation is
\begin{equation}
    I = \frac{\mathrm{PL} - \mathrm{PL}_{\mathrm{thr}}}
    {\mathrm{PL}_{\max} - \mathrm{PL}_{\mathrm{thr}}}
    = \frac{\mathrm{PL} + 104}{29}.
\end{equation}
Hence,
\begin{equation}
    \mathrm{PL} = 29I - 104.
\end{equation}
Since RMSE depends only on differences, the additive offset cancels and
\begin{equation}
    \mathrm{RMSE}_\mathrm{dB} = 29 \times \mathrm{RMSE}_\mathrm{norm}.
\end{equation}
NMSE is not directly convertible without knowledge of the signal variance.

\subsection{NLoS Difficulty and Why the Challenge Values Look Narrow}

The ICASSP 2023 challenge still contains NLoS effects; the low reported RMSE
values do not mean that NLoS is solved in general. Rather, the benchmark
compresses the hard tail relative to our setting:
\begin{enumerate}[nosep]
    \item the metric is computed in a 29~dB image-domain window,
    \item the transmitter is always rooftop-mounted, which reduces extreme
          low-altitude urban shadowing,
    \item building pixels are masked to zero error,
    \item train and test cities come from the same procedural generator.
\end{enumerate}

Therefore, NLoS is still present, but its contribution to the reported RMSE is
much less extreme than in our city-holdout UAV setting. A value around 1~dB on
this challenge should not be interpreted as evidence that difficult urban NLoS
propagation is essentially solved.

\subsection{Quantisation as a Real Limitation}

Label precision is also part of the comparison story.

For the RadioMap3DSeer challenge, the 29~dB rooftop range is stored in
\texttt{uint8}, so one gray-level step corresponds to
\[
    \frac{29}{255} = 0.1137~\mathrm{dB}
\]
per level, which is relatively fine.

For our dataset, by contrast, the ground truth is effectively quantised at about
1~dB resolution. This is substantially coarser and is a real limitation:
\begin{itemize}[nosep]
    \item the irreducible quantisation-noise floor is
    \[
        \sigma_q = \frac{\Delta}{\sqrt{12}} = \frac{1}{\sqrt{12}} = 0.2887~\mathrm{dB},
    \]
    \item sharp LoS/NLoS transitions are snapped to integer bins,
    \item very small improvements can be obscured by the discretisation of the labels.
\end{itemize}

This does not explain a 9~dB RMSE on its own, but it does impose a real ceiling
on the final achievable accuracy and makes sub-dB gains harder both to learn and
to measure reliably.

\subsection{Fair-Comparison Criteria}

The most common mistake in this literature is to compare headline RMSE values
without checking whether the underlying task is actually the same. For supervisor-
facing discussion, Table~\ref{tab:fairness} is the most useful compact summary.

\begin{table}[h!]
\centering
\caption{Fair-comparison checklist across representative methods.}
\label{tab:fairness}
\small
\begin{tabularx}{\textwidth}{Xcccccc}
\toprule
\textbf{Method} & \textbf{Geometry / labels} & \textbf{Unseen city} & \textbf{Cont. height}
& \textbf{512+ px} & \textbf{Calib.} & \textbf{Comparable} \\
\midrule
PMNet / ICASSP 2023 & Procedural + simulated & Partial & No & No & No & No \\
REM-Net+ / RMTransformer & Procedural + simulated & Partial & No & No & No & No \\
Gao et al. & Procedural + simulated & No & No & No & No & No \\
AIRMap & Measured / calibrated & Yes & No & No & Partial & Partial \\
Tarhouni et al. & Real-city maps + simulated labels & Yes & No & No & No & \better{Closest} \\
WiCo-PG & Proprietary / multimodal & Unclear & Yes & Unclear & Unclear & Partial \\
\ours{Try 66} & \ours{Real-city geometry + simulated labels} & \ours{Yes} & \ours{Yes} & \ours{Yes} & \ours{No} & \ours{Reference} \\
\bottomrule
\end{tabularx}
\end{table}

The table shows why a 1--2~dB result on a same-generator fixed-height benchmark
is not automatically stronger than a 9.41~dB result on unseen city geometries with
continuous UAV altitude. The protocol difficulty is materially different.

\subsection{Continuous UAV Height as a Separate Generalisation Axis}

Continuous UAV height conditioning is not a minor metadata extension over the
fixed-Tx literature. It changes the propagation regime itself:
\begin{itemize}[nosep]
    \item low altitude: strong blockage and long NLoS corridors,
    \item medium altitude: mixed diffraction / rooftop clearance behaviour,
    \item high altitude: smoother, increasingly LoS-dominant propagation.
\end{itemize}

Therefore, Try~66 is not solving a single radio-map distribution. It is solving
a \emph{family} of altitude-indexed distributions over the 12--478~m range.
Fixed-height benchmarks only solve one slice of that family. This is one of the
strongest novelty claims of the project and should be highlighted explicitly.

%---------------------------------------------------------------------
\section{Outdoor Benchmarks}
\label{sec:outdoor}
%---------------------------------------------------------------------

\subsection{ICASSP 2023 Outdoor Pathloss Challenge}
\label{ssec:icassp2023}

\textbf{Dataset:} RadioMap3DSeer \cite{icassp2023challenge}. Synthetic city maps
for Ankara, Berlin, Glasgow, Ljubljana, London, Tel Aviv (training); Istanbul (test).
All cities generated by the \emph{same} procedural simulator (Wireless InSite
ray-tracer). Resolution: 256$\times$256. The published metric is computed in the
normalised image domain corresponding to a 29~dB rooftop-dataset interval.
Tx height: fixed rooftop 6--20~m. Building interior pixels masked to zero error.

\begin{table}[h!]
\centering
\caption{ICASSP 2023 Outdoor Challenge results (normalised RMSE and exact dB equivalent from the RadioMap3DSeer mapping).}
\label{tab:icassp2023}
\begin{tabular}{lccl}
\toprule
\textbf{Method} & \textbf{RMSE (norm.)} & \textbf{RMSE (dB)} & \textbf{Notes} \\
\midrule
RMTransformer \cite{rmtransformer2025} & 0.0270 & \better{0.7830} & Transformer-CNN hybrid \\
REM-Net+ \cite{remnet2024} & 0.0349 & 1.0121 & Late submission, 2024 \\
PMNet \cite{pmnet2023} & 0.0383 & 1.1107 & ICASSP 2023 1st place \\
Agile/REM-U-Net & 0.0451 & 1.3079 & \\
PPNet & 0.0507 & 1.4703 & \\
Gao et al.~\cite{gao2026} & $\approx$0.030 & 0.8700 & Est. from 1.2--3.0~dB gain claim \\
\midrule
\ours{Try 66 (ours)} & \ours{---} & \ours{9.41~dB} & \ours{Different problem (see below)} \\
\bottomrule
\end{tabular}
\end{table}

\textbf{Why this is \emph{not} a valid comparison to Try 66:}

\begin{enumerate}[nosep]
    \item \textbf{Same-distribution test.} Istanbul is generated by the same
          ray-tracing simulator as training cities. The model implicitly memorises
          the generator's statistical patterns rather than learning cross-city
          generalisation. Our city-holdout still uses simulator-generated labels,
          but over OSM-derived city geometries with genuinely different morphology
          never seen during training.

    \item \textbf{Fixed Tx height.} No height generalisation is required.
          Our model must generalise across 12--478~m of continuous UAV altitude,
          which is a fundamentally different conditioning problem.

    \item \textbf{Building pixels masked.} Interior pixels are excluded from
          evaluation, reducing apparent RMSE by an estimated 15--20\%. Our metrics
          are computed over all valid ground pixels.

    \item \textbf{Narrow prediction range.} The challenge metric lives in the
          RadioMap3DSeer rooftop scaling window of 29~dB, which is much narrower
          than our 65--130~dB effective range, making RMSE values intrinsically lower.
\end{enumerate}

\subsection{Gao et al.\ --- Multi-Layer Segmentation with Weighting Map}

\textbf{Paper:} ``Effective outdoor pathloss prediction: A multi-layer segmentation
approach with weighting map''~\cite{gao2026}. arXiv:2601.08436, January 2026.

\textbf{Method:} ResNet-based deep learning with Tx/Rx depth maps and a propagation
corridor weighting map emphasising pixels along the direct Tx-Rx path.
Claims 1.2--3.0~dB improvement over PPNet, RPNet, and ViT baselines on the ICASSP
2023 and ITU Challenge 2024 datasets. Achieves 60\% FLOPs reduction.

\textbf{Relevance:} This paper directly motivated our inclusion of \texttt{tx\_depth\_map}
and corridor-based loss weighting in Try 66. Both are implemented and operational.

\subsection{PathFinder --- Multi-Transmitter Generalisation}

\textbf{Paper:} PathFinder \cite{pathfinder2025}. arXiv:2512.14150, December 2025.

\textbf{Method:} Disentangled Feature Encoding (DFE) + Mask-Guided Low-Rank Attention
(MLA). Introduces the Single-to-Multi-Transmitter (S2MT-RPP) benchmark to measure
generalisation under distribution shift (training on 1 Tx, testing on 4 Tx).

\textbf{Result:} Outperforms PMNet and PPNet on multi-Tx generalisation.
Exact dB improvement not reported; results shown via visualisation.

\textbf{Relevance:} Validates that distribution-shift generalisation is a hard
open problem. Even the best 2025 methods show artefacts in unseen configurations.

\subsection{AIRMap --- Wireless Digital Twin Generalisation}

\textbf{Paper:} AIRMap \cite{airmap2025}. arXiv:2511.05522, November 2025.

\textbf{Method:} U-Net autoencoder with 2D elevation maps only (terrain + building
heights). No LoS mask, no depth maps, no height conditioning.
Trained on 1.2M Boston-area samples; tested on 4 distinct unseen urban and
rural environments (zero-shot, no fine-tuning).

\textbf{Result:} $< 4$~dB RMSE on path gain in unseen environments. Inference
in 4~ms (NVIDIA L40S). With 20\% field measurement calibration: $\approx$5\%
median relative error.

\textbf{Assessment:} The sub-4~dB result is after applying field calibration,
which corrects systematic bias. Their zero-shot (uncalibrated) result is not
explicitly reported. Our 9.41~dB is fully zero-shot with no calibration data.
The gap is plausibly explained by: (i) our harder evaluation (absolute path loss vs.
calibrated path gain), (ii) 66 city diversity vs.\ Boston-only training, and
(iii) our continuous UAV height conditioning task.

\subsection{Tarhouni et al.\ --- Out-of-Domain Suburban Transfer}

\textbf{Paper:} \cite{tarhouni2025}. arXiv:2510.00696, October 2025.

\textbf{Method:} CNN-based path loss prediction trained on KAUST suburban campus,
tested on Rio de Janeiro (entirely unseen city). Resolution: 64$\times$64.
Fixed terrestrial Tx, 2.4~GHz.

\begin{table}[h!]
\centering
\caption{Tarhouni et al.\ cross-domain results vs.\ Try 66.}
\label{tab:tarhouni}
\begin{tabularx}{\textwidth}{lXccc}
\toprule
\textbf{Method} & \textbf{Train} & \textbf{Test} & \textbf{Res.} & \textbf{RMSE} \\
\midrule
Tarhouni et al. & KAUST (1 suburban city) & Rio de Janeiro & 64$\times$64 & 9.3--11.3~dB \\
\ours{Try 66 (ours)} & \ours{66 cities (holdout)} & \ours{12 unseen cities} & \ours{513$\times$513} & \ours{9.41~dB} \\
\bottomrule
\end{tabularx}
\end{table}

\textbf{Assessment:} This is the \emph{most honest published comparison} for
cross-city generalisation. Their result ranges 9.3--11.3~dB on a 64$\times$64 map
with a single-campus training environment; our 9.41~dB is at 513$\times$513 with
66-city training diversity and continuous UAV height conditioning. Our result is
competitive or superior across all measured dimensions.

\subsection{WiCo-PG --- UAV-to-Ground Foundation Model}

\textbf{Paper:} WiCo-PG \cite{wicopg2025}. arXiv:2511.15030, November 2025.

\textbf{Method:} Dual VQGANs + Transformer with frequency-guided Mixture of Experts.
Trained on a proprietary multi-modal dataset with UAV altitudes and diverse frequencies.
Uses RGB drone camera images as auxiliary input for cross-modal conditioning.

\textbf{Result:} NMSE~=~0.012; outperforms LLM-based baseline by $>$6.98~dB
(improvement metric, not absolute). Few-shot (2.7\% samples): $>$1.37~dB
over LLM4PG.

\textbf{Relevance:} One of the only published works that explicitly addresses
variable UAV altitude in path loss prediction. Their Mixture of Experts for
multi-frequency is conceptually related to our sinusoidal FiLM for continuous height.
Direct comparison is not possible because NMSE is reported rather than dB RMSE.

%---------------------------------------------------------------------
\section{Indoor Benchmarks}
\label{sec:indoor}
%---------------------------------------------------------------------

\subsection{ICASSP 2025 Indoor Pathloss Challenge}

\textbf{Dataset:} Multi-frequency indoor (868~MHz, 1.8~GHz, 3.5~GHz), directive
antennas, 3 tasks of increasing spatial complexity
\cite{icassp2025indoor,transPathNet2025,ippnet2025}.

\begin{table}[h!]
\centering
\caption{ICASSP 2025 Indoor Challenge results (weighted average across 3 tasks).}
\label{tab:icassp2025}
\begin{tabular}{lcc}
\toprule
\textbf{Method} & \textbf{RMSE (dB, wtd.\ avg)} & \textbf{Task 1 / 2 / 3} \\
\midrule
Wenlihan Lu (HKUST-GZ) & \better{9.411} & 7.85 / 10.07 / 10.09 \\
IPP-Net \cite{ippnet2025} & 9.501 & \\
TerRaIn & 10.325 & \\
TransPathNet \cite{transPathNet2025} & 10.40 & \\
\midrule
\ours{Try 66 -- open\_sparse\_lowrise} & \ours{9.41} & \ours{(single expert, outdoor)} \\
\bottomrule
\end{tabular}
\end{table}

The numerical coincidence is notable: Try 66 achieves \textbf{9.41~dB}
(epoch 96), exactly matching the indoor challenge winner, despite:
\begin{itemize}[nosep]
    \item Operating outdoors at 513$\times$513 (vs.\ indoor, smaller maps)
    \item Requiring city-level generalisation (vs.\ fixed building layout)
    \item Conditioning on continuous UAV altitude (vs.\ fixed Tx placement)
    \item Using zero calibration data (vs.\ challenge-matched test distribution)
\end{itemize}

%---------------------------------------------------------------------
\section{Novelty of Try 66 in the Literature}
\label{sec:novelty}
%---------------------------------------------------------------------

Table~\ref{tab:novelty} summarises the key properties of Try 66 against all
relevant methods. No published paper as of April 2026 simultaneously satisfies
all six conditions.

\begin{table}[h!]
\centering
\caption{Feature comparison: Try 66 vs.\ published methods.
\yes~= satisfied, \no~= not satisfied.}
\label{tab:novelty}
\small
\begin{tabular}{lcccccc}
\toprule
\textbf{Feature} & \textbf{PMNet} & \textbf{AIRMap} & \textbf{WiCo-PG}
    & \textbf{Tarhouni} & \textbf{Gao} & \ours{\textbf{Try 66}} \\
\midrule
Outdoor (not indoor) & \yes & \yes & \yes & \yes & \yes & \ours{\yes} \\
$\geq$512$\times$512 px & \no & \no & \no & \no & \no & \ours{\yes} \\
Continuous UAV height & \no & \no & partial & \no & \no & \ours{\yes} \\
City-level holdout & \no & partial & \no & \yes & \no & \ours{\yes} \\
$\geq$20 training cities & \no & \no & \no & \no & \no & \ours{\yes} \\
No test-time calibration & \yes & \no & n/a & \yes & \yes & \ours{\yes} \\
\bottomrule
\end{tabular}
\end{table}

The sinusoidal positional encoding applied to Tx height and fed into every
FiLM layer is, to the best of our knowledge, not described in any prior work
on radio map prediction. WiCo-PG uses separate expert networks per altitude bin
(discrete); RadioLAM \cite{radiolam2025} uses discrete height bins as separate
generation targets; our approach is the only one treating height as a continuous
parameter that modulates internal feature statistics throughout the network.

%---------------------------------------------------------------------
\section{Internal Progress}
\label{sec:progress}
%---------------------------------------------------------------------

\begin{table}[h!]
\centering
\caption{Internal benchmark across training iterations (513$\times$513 unless noted).}
\label{tab:progress}
\begin{tabular}{lllcccc}
\toprule
\textbf{Try} & \textbf{Expert} & \textbf{Res.} & \textbf{Overall} & \textbf{LoS} & \textbf{NLoS} & \textbf{Notes} \\
\midrule
22 & global & 256$\times$256 & 19.94~dB & 3.78~dB & 34.43~dB & U-Net baseline \\
42 & global & varies & 19.78~dB & 3.86~dB & 34.47~dB & PMNet + prior \\
55 & lowrise & varies & 10.53~dB & 3.76~dB & 35.82~dB & PMHHNet expert \\
64 & lowrise & \textit{128$\times$128} & \textit{7.76~dB} & \textit{2.97~dB} & \textit{24.91~dB} & Coarse-only$^*$ \\
65 & lowrise & 256$\times$256 & 11.36~dB & 4.62~dB & 37.88~dB & Grokking test \\
\midrule
\ours{66} & \ours{lowrise} & \ours{513$\times$513} & \ours{9.41~dB} & \ours{3.75~dB} & \ours{31.48~dB} & \ours{Current best} \\
66 & highrise & 513$\times$513 & 28.8~dB & 4.10~dB & 34.72~dB & Epoch 98, early \\
\bottomrule
\end{tabular}
\vspace{0.3em}
\footnotesize $^*$Try 64 at 128$\times$128 is not comparable: building-edge NLoS pixels
are spatially blurred at coarse resolution, artificially reducing RMSE.
\end{table}

\subsection{Why the Remaining Gap is Mostly NLoS}

The missing 4--5~dB to the 5~dB target is not due to uniformly poor predictions
across the map. It is concentrated in a relatively small subset of structurally
hard pixels:
\begin{itemize}[nosep]
    \item LoS regions are already reasonably well modelled,
    \item the physics prior is helpful in open areas,
    \item the dominant residual error occurs near building corners, street canyons,
          and deep shadow regions where diffraction is the governing mechanism.
\end{itemize}

This is why the target is simultaneously plausible and difficult. The model does
not need a broad improvement everywhere; it needs to solve the most physics-
sensitive urban shadowing cases that dominate the quadratic loss.

\subsection{Canonical Failure Modes}

\begin{table}[h!]
\centering
\caption{Three canonical failure modes that should be discussed explicitly.}
\label{tab:failuremodes}
\begin{tabularx}{\textwidth}{lXX}
\toprule
\textbf{Failure mode} & \textbf{Typical symptom} & \textbf{Interpretation} \\
\midrule
Street canyon shadowing & Prediction remains too optimistic along long blocked corridors & Missing explicit multi-edge diffraction prior \\
Building-corner transition & Sharp local error at the LoS$\rightarrow$NLoS boundary & Boundary physics is under-modelled \\
Deep block shadow & Systematic underestimation behind dense highrise clusters & Prior mismatch dominates residual learning \\
\bottomrule
\end{tabularx}
\end{table}

These examples are more persuasive than a generic statement such as ``NLoS is
hard'', because they connect the error directly to known propagation mechanisms.

%---------------------------------------------------------------------
\section{Emerging Techniques for Closing the Gap}
\label{sec:emerging}
%---------------------------------------------------------------------

The remaining 5.4~dB gap between our 9.41~dB and the theoretical floor
($\approx$3--4~dB for lowrise) is dominated by NLoS error.
Table~\ref{tab:emerging} summarises six recent techniques that directly
target this bottleneck.

\begin{table}[h!]
\centering
\caption{Emerging research directions ranked by expected impact.}
\label{tab:emerging}
\small
\begin{tabularx}{\textwidth}{clXcc}
\toprule
\textbf{\#} & \textbf{Technique} & \textbf{Key paper} & \textbf{Gain (est.)} & \textbf{Complexity} \\
\midrule
1 & Knife-edge diffraction prior & Geometry-assisted DL \cite{geomDL2024} & 1--3~dB NLoS & High \\
2 & PDE residual loss (PINN) & ReVeal \cite{reveal2025} & 0.5--2~dB NLoS & Medium \\
3 & Dual LoS / NLoS head & (synthesised) & 1--2~dB NLoS & Medium \\
4 & Test-time adaptation & RadioPiT \cite{radiopit2025} & 0.3--1~dB all & Low--Med \\
5 & Diffusion refinement & RadioDiff-k$^2$ \cite{radiodiff2025} & 0.5--2~dB all & High \\
6 & Foundation pre-training & FM-RME \cite{fmrme2026} & 0.3--1~dB all & High \\
\bottomrule
\end{tabularx}
\end{table}

\paragraph{Knife-edge diffraction features (\#1).}
Our current physics prior is free-space path loss with corrections, containing
no diffraction information. A multi-screen knife-edge model (ITU-R P.526) can
pre-compute per-pixel diffraction loss from building geometry and provide it as
1--2 additional input channels. This directly addresses the root cause of high
NLoS error: the model must currently learn diffraction entirely from data.
The geometry-assisted approach of \cite{geomDL2024} demonstrated 10--18\%
accuracy improvement and 50\% fewer fine-tuning epochs with this technique.

\paragraph{Physics-informed PDE loss (\#2).}
ReVeal~\cite{reveal2025} derives a second-order PDE for the signal field and
adds its residual as regularisation. With only 30 training points over 514~km$^2$,
it achieves 1.95~dB RMSE in rural/suburban environments. Integrating a similar
Helmholtz-inspired term into our loss could smooth NLoS transition zones where
the model currently produces discontinuous predictions.

\paragraph{Test-time adaptation (\#4).}
RadioPiT~\cite{radiopit2025} achieves a 21.9\% RMSE reduction via test-time
adaptation (TTA). For Try~66, TTA could exploit the known accuracy of the
physics prior in LoS regions: disagreement between the model and prior in LoS
pixels signals feature miscalibration for the current city. Adapting only the
FiLM / batch normalisation parameters (1--5 gradient steps) at inference time
could improve unseen-city performance at near-zero cost.

%---------------------------------------------------------------------
\section{Honest Assessment}
\label{sec:assessment}
%---------------------------------------------------------------------

\begin{table}[h!]
\centering
\caption{Positioning of Try 66 against key benchmarks.}
\label{tab:assessment}
\begin{tabularx}{\textwidth}{Xccl}
\toprule
\textbf{Benchmark} & \textbf{Their RMSE} & \textbf{Our RMSE} & \textbf{Verdict} \\
\midrule
ICASSP 2023 (same-dist.) & 0.7830--1.1107~dB & 9.41~dB & Different problem --- invalid \\
AIRMap (calibrated unseen) & $<$4~dB & 9.41~dB & Different protocol --- not fair \\
Tarhouni (unseen city, 64$\times$64) & 9.3--11.3~dB & 9.41~dB & \better{Competitive or better} \\
ICASSP 2025 indoor winner & 9.411~dB & 9.41~dB & \better{Numerically tied, harder task} \\
WiCo-PG (UAV, NMSE) & NMSE 0.012 & --- & Incomparable metric \\
\bottomrule
\end{tabularx}
\end{table}

\textbf{Conclusion:} Try 66 is state-of-the-art for its specific problem formulation.
No published method simultaneously achieves sub-10~dB RMSE on (1) outdoor environments
at (2) full 512+ pixel resolution with (3) continuous UAV height and (4) city-level
holdout over unseen city geometries. The 5~dB overall RMSE target is challenging but physically
achievable through improvements in NLoS prior modelling and extended training.

%---------------------------------------------------------------------
\section*{References}
\label{sec:references}
%---------------------------------------------------------------------

\begin{thebibliography}{99}

\bibitem{pmnet2023}
J.~Lee, A.~F.~Molisch,
``PMNet: Robust pathloss map prediction via supervised learning,''
\textit{Proc. IEEE ICASSP 2023}, pp.~1--5.
arXiv:2211.10527.

\bibitem{remnet2024}
REM-Net+, Authorea TechRxiv, 2024.
\url{https://www.authorea.com/doi/full/10.36227/techrxiv.172226872.27268975}

\bibitem{rmtransformer2025}
RMTransformer, 2025.
arXiv:2501.05190.

\bibitem{icassp2023challenge}
C.~Yapar \textit{et al.},
``First pathloss radio map prediction challenge,''
2023.
arXiv:2310.07658.

\bibitem{dataset2212}
Ç.~Yapar, R.~Levie, G.~Kutyniok, and G.~Caire,
``Dataset of Pathloss and ToA Radio Maps With Localization Application,''
2022.
arXiv:2212.11777.

\bibitem{icassp2025indoor}
S.~Bakirtzis \textit{et al.},
``ICASSP 2025 indoor radio map challenge,''
2025.
arXiv:2501.13698.

\bibitem{ippnet2025}
IPP-Net,
\textit{IEEE ICASSP 2025}, 2nd place indoor challenge.
IEEE Xplore: 10890663.

\bibitem{transPathNet2025}
TransPathNet,
2025.
arXiv:2501.16023.

\bibitem{gao2026}
Y.~Gao \textit{et al.},
``Effective outdoor pathloss prediction: A multi-layer segmentation approach
with weighting map,''
2026.
arXiv:2601.08436.

\bibitem{tarhouni2025}
M.~Tarhouni \textit{et al.},
``Out-of-domain path loss prediction,''
2025.
arXiv:2510.00696.

\bibitem{airmap2025}
AIRMap: AI-Generated Radio Maps for Wireless Digital Twins,
2025.
arXiv:2511.05522.

\bibitem{pathfinder2025}
PathFinder: Advancing Path Loss Prediction for Single-to-Multi-Transmitter Scenario,
2025.
arXiv:2512.14150.

\bibitem{wicopg2025}
WiCo-PG: Wireless Channel Foundation Model for Pathloss Map Generation,
2025.
arXiv:2511.15030.

\bibitem{radiolam2025}
RadioLAM: MoE Diffusion for 3D Radio Maps,
2025.
arXiv:2509.11571.

\bibitem{leemolisch2023}
J.~Lee, A.~F.~Molisch,
``A scalable and generalizable pathloss map prediction,''
2023.
arXiv:2312.03950.

\bibitem{ckmimagenet2025}
CKMImageNet dataset (height metadata),
2025.
arXiv:2504.09849.

\bibitem{reveal2025}
ReVeal: A Physics-Informed Neural Network for High-Fidelity Radio Environment Mapping,
2025.
arXiv:2502.19646.

\bibitem{geomDL2024}
Diffraction and Scattering Aware Radio Map and Environment Reconstruction
using Geometry Model-Assisted Deep Learning,
\textit{IEEE TWC}, 2024.
arXiv:2403.00229.

\bibitem{radiodiff2025}
RadioDiff: Conditional Denoising Diffusion for Radio Map Construction,
\textit{IEEE TCCN / JSAC}, 2025--2026.
Code: \url{https://github.com/UNIC-Lab/RadioDiff}.

\bibitem{radiopit2025}
RadioPiT: Radio Map Generation with Pixel Transformer Driven by Ultra-Sparse
Real-World Data,
2025.
arXiv:2512.01451.

\bibitem{fmrme2026}
FM-RME: Foundation Model Empowered Radio Map Estimation,
2026.
arXiv:2602.22231.

\bibitem{rmgen2025}
RM-Gen: Denoising Diffusion Probabilistic Model for Radio Map Estimation,
2025.
arXiv:2501.06604.

\end{thebibliography}

\end{document}
